Choosing a pre-trained model for a new machine-learning problem is, in practice, a manual and largely unprincipled search: a practitioner picks a handful of familiar models, fine-tunes each on their data, and keeps whichever wins. This does not scale to a model zoo of millions of models, and it ignores real-world constraints such as parameter budget, license, or modality that should prune the search space before any compute is spent.

Recent work reframes model selection as \emph{link prediction} on a heterogeneous model--dataset graph: TransferGraph \citep{li2024transfergraph} trains a regressor over Node2Vec embeddings of a model--dataset graph to predict fine-tuning accuracy, and ArtifactLinker \citep{yu2026artifactlinker} trains a joint graph neural network (GNN) with a link-prediction head and an attribute-regression head over an artifact graph of 14{,}053 models, datasets, papers, and code repositories (51{,}337 relations), scoring every unobserved (model, dataset) pair as \emph{link probability} $\times$ \emph{predicted attribute value}. Two limitations remain. First, both methods only work for datasets already present as nodes: a genuinely new, cold-start dataset has no edges and is unscorable. Second, both stop at a predicted score; nothing in the pipeline confirms that the top-ranked model actually performs well when it is fine-tuned or evaluated on the target data.

We present \textbf{AutoModel Advisor}, an agentic system that addresses both gaps. Given a user's dataset, task, and constraints (e.g., a maximum parameter count), an orchestrator coordinates two agents. The \textbf{Retrieval Agent} makes a cold-start dataset scorable by summarizing it with an LLM, embedding the summary, and either (i) retrieving its top-$K$ most similar dataset neighbors already in ArtifactLinker's graph and prompting an LLM with their measured model performance (\emph{cosine retrieval}), or (ii) inserting the dataset as a new node into the graph and letting ArtifactLinker's trained GNN score every candidate model inductively (\emph{direct GNN scoring}). The \textbf{Evaluation Agent} then writes and runs real evaluation code for the shortlisted candidates on the user's data, replacing a predicted score with a measured one. Critically, the loop does not stop after one pass: the orchestrator feeds the Evaluation Agent's per-model results and error analysis (which slices, classes, or examples a candidate fails on) back to the Retrieval Agent, which uses that evidence to confirm, drop, or replace candidates in the next round.

Our contributions are: (1) a two-agent, orchestrated architecture that turns single-shot graph scoring into closed-loop, de-facto-verified model selection; (2) two complementary cold-start strategies that make an out-of-graph dataset scorable by ArtifactLinker's inductive machinery; (3) an evaluation protocol that isolates the contribution of each evidence source (dataset-similarity retrieval vs. GNN prediction) and compares the full agentic system against a raw-GNN baseline and a strong general-purpose coding-agent baseline under matched compute and evaluation-protocol constraints.